\emph{प्रथम-क्रम भाषेत} \emph{स्थिर}, \emph{फलन} आणि \emph{विधेय} चिन्हे असतात. प्रत्येक फलनचिन्हाची आणि विधेयचिन्हाची स्थानसंख्या ठरलेली असते. उदाहरणार्थ, \emph{अंकगणिताच्या भाषेत} एकच स्थिरचिन्ह~$\Obj 0$, एक एकस्थानी फलनचिन्ह $\prime$, दोन द्विस्थानी फलनचिन्हे~$+$ आणि $\times$, आणि एक द्विस्थानी विधेयचिन्ह~$<$ आहे. \emph{चर}, स्थिरचिन्हे आणि फलनचिन्हे यांपासून आपण भाषेची \emph{पदे} बनवतो. भाषेची पदे, विधेयचिन्हे आणि \emph{एकरूपता चिन्ह}~$\eq$ यांच्या साहाय्याने \emph{आण्विक सूत्रे} बनवतो. त्या सूत्रांपासून तार्किक संयोजक $\lnot$, $\lor$, $\land$, $\lif$, $\liff$ आणि संख्यापक $\lforall$ व $\lexists$ वापरून भाषेची सूत्रे बनवतो. पदे आणि सूत्रे तयार करताना आवश्यक कंस नेहमी घातले जात असल्यामुळे प्रत्येक सूत्र वाचण्याची नेमकी एकच रीत असते. त्यामुळे सूत्रांच्या रचनेवरील विगमनाने व्याख्या देता येतात.

सूत्रांतील चरांच्या आस्थिती कधी कधी संबंधित संख्यापकाच्या अधीन असतात. एखाद्या चराची आस्थिती संख्यापकाच्या \emph{व्याप्तीत} येत असेल, तर ती \emph{बद्ध} मानली जाते, अन्यथा \emph{मुक्त} मानली जाते. या सर्व संकल्पनांच्या व्याख्या विगमनाधारित आहेत. सूत्रात चराच्या जागी पदाचे \emph{आदेशन} करण्याची क्रियाही आपण विगमनाने परिभाषित करतो. मुक्त चराच्या आस्थिती नसलेल्या सूत्रांना \emph{वाक्ये} म्हणतात.
